% Part: turing-machines
% Chapter: undecidability
% Section: verification

\documentclass[../../../include/open-logic-section]{subfiles}

\begin{document}

\olfileid{tur}{und}{ver}
\olsection{بررسی بازنمایی}

\begin{explain}
برای بررسی اینکه بازنمایی ما کار می‌کند باید دو چیز را ثابت کنیم.
نخست باید نشان دهیم اگر $M$ روی ورودی~$w$ متوقف شود، آنگاه
$!T(M,w)\lif !E(M,w)$ معتبر است. سپس باید عکس آن را نشان دهیم؛ یعنی
اگر $!T(M,w)\lif !E(M,w)$ معتبر باشد، $M$ در اجرای روی ورودی~$w$ در
واقع سرانجام متوقف می‌شود.

راهبرد اثبات این دو نتیجه بسیار متفاوت است. برای نتیجهٔ نخست باید
نشان دهیم !!a{sentence}ای از منطق مرتبه‌اول (یعنی
$!T(M,w)\lif !E(M,w)$) معتبر است. ساده‌ترین راه، ارائهٔ !!a{derivation}ی
برای آن است. اما قرار است اثبات ما برای همهٔ $M$ و~$w$ کار کند؛ پس
در واقع یک !!{sentence}ٔ منفرد نداریم که برایش !!a{derivation} بدهیم، بلکه
بی‌نهایت جمله داریم. بنابراین بهترین کاری که می‌توانیم بکنیم این است
که با استقرا ثابت کنیم $M$ و~$w$ هرچه باشند و توقف~$M$ روی ورودی~$w$
هرچند گام طول بکشد، !!a{derivation}ی از $!T(M,w)\lif !E(M,w)$ وجود دارد.

طبعاً استقرای ما بر تعداد گام‌هایی پیش می‌رود که $M$ تا رسیدن به
پیکربندی توقف برمی‌دارد. در اثبات استقرایی نشان می‌دهیم به ازای هر
گام~$n$ از اجرای~$M$ روی ورودی~$w$ تا و شامل نخستین پیکربندی توقف،
$!T(M,w)\Entails !C(M,w,n)$، که در آن $!C(M,w,n)$ پیکربندی~$M$ را پس
از اجرای~$n$ گام روی~$w$ به‌درستی توصیف می‌کند. حال اگر $M$ مثلاً پس
از~$n$ گام روی~$w$ متوقف شود، $!C(M,w,n)$ پیکربندی توقف را توصیف
خواهد کرد. همچنین نشان می‌دهیم هرگاه $!C(M,w,n)$ پیکربندی توقف را
توصیف کند، $!C(M,w,n)\Entails !E(M,w)$. پس اگر $M$ روی ورودی~$w$
متوقف شود، برای بعضی~$n$ پس از~$n$ گام در پیکربندی توقف خواهد بود.
بنابراین $!T(M,w)\Entails !C(M,w,n)$، که در آن $!C(M,w,n)$ پیکربندی
توقف را توصیف می‌کند؛ و چون در این حالت
$!C(M,w,n)\Entails !E(M,w)$، به
$!T(M,w)\Entails !E(M,w)$، یعنی
$\Entails !T(M,w)\lif !E(M,w)$، می‌رسیم.

راهبرد اثبات عکس بسیار متفاوت است. اینجا فرض می‌کنیم
$\Entails !T(M,w)\lif !E(M,w)$ و باید ثابت کنیم $M$ روی ورودی~$w$
متوقف می‌شود. از فرض نتیجه می‌گیریم
$!T(M,w)\Entails !E(M,w)$؛ یعنی $!E(M,w)$ در هر !!{structure}ی که
$!T(M,w)$ در آن صادق است، صادق خواهد بود. پس !!a{structure}ی چون
$\Struct{M}$ توصیف می‌کنیم که $!T(M,w)$ در آن صادق است: دامنهٔ آن
$\Nat$ خواهد بود و تعبیر همهٔ $\Obj Q_q$ها و~$\Obj S_\sigma$ها را
پیکربندی‌های~$M$ هنگام اجرای روی ورودی~$w$ تعیین می‌کنند. برای نمونه،
$\Sat{M}{\Obj Q_q(\num{m},\num{n})}$ اگر و تنها اگر $M$ پس از اجرای
$n$ گام روی ورودی~$w$ در حالت~$q$ باشد و خانهٔ~$m$ را پویش کند. اکنون
چون بنا بر فرض $!T(M,w)\Entails !E(M,w)$، و بنا بر ساخت
$\Sat{M}{!T(M,w)}$، داریم $\Sat{M}{!E(M,w)}$. اما
$\Sat{M}{!E(M,w)}$ اگر و تنها اگر $n\in\Domain{M}=\Nat$ای وجود داشته
باشد که $M$ پس از اجرای~$n$ گام روی ورودی~$w$ در پیکربندی توقف باشد.
\end{explain}

\begin{defn}
بگذارید $!C(M,w,n)$ !!{sentence}ٔ زیر باشد:
\[
\Obj Q_q(\num{m}, \num{n}) \land \Obj S_{\sigma_0}(\num{0}, \num{n})
\land \dots \land \Obj S_{\sigma_k}(\num{k}, \num{n}) \land
\lforall[x][(\num{k} < x \lif \Obj S_\TMblank(x, \num{n}))]
\]
که در آن $q$ حالت~$M$ در زمان~$n$ است، $M$ در زمان~$n$ خانهٔ~$m$ را
پویش می‌کند، خانهٔ~$i$ در زمان~$n$ برای $0\le i\le k$ شامل نماد
$\sigma_i$ است، و~$k$ بزرگ‌ترِ این دو است: شمارهٔ راست‌ترین خانهٔ
ناخالیِ نوار در زمان~$0$ و شمارهٔ راست‌ترین خانه‌ای که سر نوار پس از
$n$ گام دیده است.
\end{defn}

\begin{lem}\ollabel{lem:halt-config-implies-halt}
اگر اجرای~$M$ روی ورودی~$w$ پس از~$n$ گام در پیکربندی توقف باشد،
$!C(M,w,n)\Entails !E(M,w)$.
\end{lem}

\begin{proof}
فرض کنید $M$ پس از~$n$ گام برای ورودی~$w$ متوقف شود. حالت~$q$،
خانهٔ~$m$ و نماد~$\sigma$ای وجود دارند به‌گونه‌ای که:
\begin{enumerate}
\item پس از~$n$ گام، $M$ در حالت~$q$ خانهٔ~$m$ را که روی آن
  $\sigma$ قرار دارد پویش می‌کند؛
\item تابع گذار~$\delta(q,\sigma)$ تعریف‌نشده است.
\end{enumerate}
$!C(M,w,n)$ توصیف این پیکربندی است و بندهای
$\Obj Q_q(\num{m},\num{n})$ و
$\Obj S_\sigma(\num{m},\num{n})$ را در بر دارد. این دو بند در کنار هم
$!E(M,w)$ را نتیجه می‌دهند:
\[
\lexists[x][\lexists[y][(\bigvee_{\tuple{q, \sigma} \in
      X}(\Obj Q_q(x, y) \land \Obj S_\sigma(x, y)))]]
\]
زیرا
$\Obj Q_q(\num{m},\num{n})\land\Obj S_\sigma(\num{m},\num{n})
\Entails \bigvee_{\tuple{q',\sigma'}\in X}
(\Obj Q_{q'}(\num{m},\num{n})\land
\Obj S_{\sigma'}(\num{m},\num{n}))$، چون
$\tuple{q,\sigma}\in X$.
\end{proof}

\begin{explain}
پس اگر $M$ برای ورودی~$w$ متوقف شود، $n$ای وجود دارد که
$!C(M,w,n)\Entails !E(M,w)$. اکنون نشان می‌دهیم برای هر زمان~$n$ که
$M$ پیش از انجام~$n$ گام متوقف نشده باشد،
$!T(M,w)\Entails !C(M,w,n)$.
\end{explain}

\begin{lem}
\ollabel{lem:config}
برای هر~$n$، اگر $M$ پیش از انجام~$n$ گام متوقف نشده باشد، آنگاه
$!T(M,w)\Entails !C(M,w,n)$.
\end{lem}

\begin{proof}
پایهٔ استقرا: اگر $n=0$، عطف‌های~$!C(M,w,0)$ عطف‌های~$!T(M,w)$ نیز
هستند و بنابراین از آن نتیجه می‌شوند.

فرض استقرا: اگر $M$ پیش از گام $n$اُم متوقف نشده باشد، آنگاه
$!T(M,w)\Entails !C(M,w,n)$. باید نشان دهیم (مگر آنکه
$!C(M,w,n)$ پیکربندی توقف را توصیف کند)
$!T(M,w)\Entails !C(M,w,n+1)$.

فرض کنید $n\ge 0$ و~$M$ که روی~$w$ آغاز شده است پس از~$n$ گام در
حالت~$q$ خانهٔ~$m$ را پویش کند. چون $M$ پس از~$n$ گام متوقف نمی‌شود،
دستوری از یکی از سه صورت زیر باید در برنامهٔ~$M$ باشد:

\begin{enumerate}
\item \ollabel{right} $\delta(q, \sigma) = \tuple{q', \sigma', \TMright}$

\item \ollabel{left} $\delta(q, \sigma) = \tuple{q', \sigma', \TMleft}$

\item \ollabel{stay} $\delta(q, \sigma) = \tuple{q', \sigma', \TMstay}$
\end{enumerate}

هر یک از این سه حالت را به‌ترتیب بررسی می‌کنیم.

\begin{enumerate}
\item فرض کنید دستوری از صورت~\olref{right} وجود دارد. بنا بر
  \olref[rep]{defn:tm-descr}\olref[rep]{rep-right}، این یعنی
\begin{align*}
& \lforall[x][\lforall[y][((\Obj Q_{q}(x,
  y) \land \Obj S_\sigma(x, y)) \lif {}]]\\
& \qquad (\Obj Q_{q'}(x',
  y') \land \Obj S_{\sigma'}(x, y') \land !A(x, y)))
\intertext{عطفی از $!T(M,w)$ است. این فرمول !!{sentence}ٔ زیر را نتیجه
  می‌دهد (با نمونه‌سازی همگانیِ $\num{m}$ برای~$x$ و~$\num{n}$
  برای~$y$):}
& (\Obj Q_{q}(\num{m}, \num{n}) \land \Obj S_{\sigma}(\num{m},
\num{n})) \lif {}\\
& \qquad (\Obj Q_{q'}(\num{m}', \num{n}') \land
\Obj S_{\sigma'}(\num{m}, \num{n}') \land !A(\num{m}, \num{n})).
\intertext{بنا بر فرض استقرا، $!T(M,w)\Entails !C(M,w,n)$؛ یعنی}
& \Obj Q_q(\num{m}, \num{n}) \land \Obj S_{\sigma_0}(\num{0}, \num{n})
\land \dots \land \Obj S_{\sigma_k}(\num{k}, \num{n}) \land \\
& \qquad\lforall[x][(\num{k} < x \lif \Obj S_\TMblank(x, \num{n}))]\\
\intertext{چون پس از~$n$ گام خانهٔ~$m$ شامل~$\sigma$ است، عطف متناظر
  $\Obj S_\sigma(\num{m},\num{n})$ است؛ پس این فرمول نتیجه می‌دهد:}
& \Obj Q_{q}(\num{m}, \num{n}) \land \Obj S_{\sigma}(\num{m},
   \num{n})
\intertext{اکنون به فرمول زیر می‌رسیم:}
& \Obj Q_{q'}(\num{m}', \num{n}') \land \Obj S_{\sigma'}(\num{m},
  \num{n}') \land {}\\
& \qquad\Obj S_{\sigma_0}(\num{0}, \num{n}') \land \dots \land
  \Obj S_{\sigma_k}(\num{k}, \num{n}') \land {}\\
& \qquad \lforall[x][(\num{k} < x \lif \Obj S_\TMblank(x, \num{n}'))]
\end{align*}
به این ترتیب: سطر نخست با وضع مقدم مستقیماً از تالیِ شرطی پیشین به
دست می‌آید. هر عطف در سطر میانی، که
$S_{\sigma_m}(\num{m},\num{n}')$ را در بر نمی‌گیرد، از عطف متناظر در
$!C(M,w,n)$ همراه با~$!A(\num{m},\num{n})$ نتیجه می‌شود.

اگر $m<k$، آنگاه
$!T(M,w)\Proves\num{m}<\num{k}$
(\olref[rep]{prop:mlessk}) و از تعدیِ~$<$ داریم
$\lforall[x][(\num{k}<x\lif\num{m}<x)]$. اگر $m=k$، همین فرمول فقط
با منطق نتیجه می‌شود. سپس سطر آخر از عطف متناظر در~$!C(M,w,n)$،
$\lforall[x][(\num{k}<x\lif\num{m}<x)]$ و~$!A(\num{m},\num{n})$
نتیجه می‌شود. اگر $m<k$، این همان~$!C(M,w,n+1)$ است.

اکنون فرض کنید $m=k$. در این صورت پس از~$n+1$ گام سر نوار خانهٔ
$k+1$ را نیز دیده است که اکنون راست‌ترین خانهٔ دیده‌شده است. پس
$!C(M,w,n+1)$ عطف تازهٔ
$\Obj S_\TMblank(\num{k}',\num{n}')$ را دارد و آخرین عطف آن
$\lforall[x][(\num{k}'<x\lif\Obj S_\TMblank(x,\num{n}'))]$ است. باید
بررسی کنیم که این دو !!{sentence}s نیز نتیجه می‌شوند.

از پیش داریم
$\lforall[x][(\num{k}<x\lif\Obj S_\TMblank(x,\num{n}'))]$. به‌ویژه،
از این
$\num{k}<\num{k}'\lif\Obj S_\TMblank(\num{k}',\num{n}')$ به دست
می‌آید. از اصل موضوعهٔ~$\lforall[x][x<x']$ نیز
$\num{k}<\num{k}'$ را می‌گیریم. با وضع مقدم،
$\Obj S_\TMblank(\num{k}',\num{n}')$ نتیجه می‌شود.

همچنین چون $!T(M,w)\Proves\num{k}<\num{k}'$، اصل موضوعهٔ تعدیِ~$<$
$\lforall[x][(\num{k}'<x\lif\Obj S_\TMblank(x,\num{n}'))]$ را به ما
می‌دهد. (بررسی این نکته را به‌عنوان تمرین وا می‌گذاریم.)

\item فرض کنید دستوری از صورت~\olref{left} وجود دارد. آنگاه بنا بر
  \olref[rep]{defn:tm-descr}\olref[rep]{rep-left}،
\begin{align*}
& \lforall[x][\lforall[y][((\Obj Q_{q}(x', y) \land \Obj
    S_{\sigma}(x', y)) \lif {}]]\\
& \qquad (\Obj Q_{q'}(x, y') \land \Obj
  S_{\sigma'}(x', y') \land !A(x', y))) \land {}\\
& \lforall[y][((\Obj Q_{q_i}(\num{0}, y) \land \Obj S_{\sigma}(\num{0},
    y)) \lif {}]\\
& \qquad (\Obj Q_{q_j}(\num{0}, y') \land \Obj S_{\sigma'}(\num{0}, y')
  \land !A(\num{0}, y)))
\intertext{عطفی از $!T(M,w)$ است. اگر $m>0$، بگذارید $l=m-1$
  (یعنی $m=l+1$). عطف نخستِ !!{sentence}ٔ بالا فرمول زیر را نتیجه می‌دهد:}
& (\Obj Q_{q}(\num{l}', \num{n}) \land \Obj S_{\sigma}(\num{l}', \num{n}))
\lif {} \\
& \qquad (\Obj Q_{q'}(\num{l}, \num{n}') \land \Obj S_{\sigma'}(\num{l}',
\num{n}') \land !A(\num{l}', \num{n}))
\intertext{در غیر این صورت بگذارید $l=m=0$ و !!{sentence}ٔ زیر را که از
  عطف دوم نتیجه می‌شود در نظر بگیرید:}
& ((\Obj Q_{q_i}(\num{0}, \num{n}) \land \Obj S_{\sigma}(\num{0}, \num{n})) \lif {}\\
& \qquad   (\Obj Q_{q_j}(\num{0}, \num{n}') \land \Obj S_{\sigma'}(\num{0}, \num{n}') \land
!A(\num{0}, \num{n})))
\intertext{هر یک از این دو جمله، مانند پیش، فرمول زیر را نتیجه می‌دهد:}
&  \Obj Q_{q'}(\num{l}, \num{n}') \land \Obj S_{\sigma'}(\num{m},
  \num{n}') \land {}\\
&\qquad  \Obj S_{\sigma_0}(\num{0}, \num{n}')\land \dots \land
  \Obj S_{\sigma_k}(\num{k}, \num{n}')  \land {} \\
& \qquad  \lforall[x][(\num{k} < x
  \lif \Obj S_\TMblank(x, \num{n}'))]
\end{align*}
(توجه کنید در حالت نخست
$\num{l}'\ident\num{l+1}\ident\num{m}$ و در حالت دوم
$\num{l}\ident\num{0}$.) اما این دقیقاً $!C(M,w,n+1)$ است.

\item حالت~\olref{stay} به‌عنوان تمرین واگذار می‌شود.
\end{enumerate}
نشان دادیم هرگاه $M$ پیش از انجام~$n$ گام متوقف نشده باشد،
$!T(M,w)\Entails !C(M,w,n)$.
\end{proof}

\begin{prob}
حالت~\olref[tur][und][ver]{stay} از اثبات
\olref[tur][und][ver]{lem:config} را کامل کنید.
\end{prob}

\begin{prob}
!!a{derivation}ی از~$\Obj S_{\sigma_i}(\num{i},\num{n}')$ بر پایهٔ
$\Obj S_{\sigma_i}(\num{i},\num{n})$ و~$!A(m,n)$ ارائه کنید (با فرض
$i\neq m$، یعنی یا $i<m$ یا $m<i$).
\end{prob}

\begin{prob}
!!a{derivation}ی از
$\lforall[x][(\num{k}' < x \lif \Obj S_\TMblank(x, \num{n}'))]$ بر
پایهٔ $\lforall[x][(\num{k} < x \lif \Obj
  S_\TMblank(x, \num{n}'))]$، $\lforall[x][x < x']$ و
$\lforall[x][\lforall[y][\lforall[z][ ((x < y \land y < z) \lif x <
      z)]]]$ ارائه کنید.)
\end{prob}


\begin{lem}
\ollabel{lem:valid-if-halt}
اگر $M$ روی ورودی~$w$ متوقف شود، آنگاه $!T(M,w)\lif !E(M,w)$ معتبر
است.
\end{lem}

\begin{proof}
از \olref{lem:config} می‌دانیم برای هر زمان~$n$ که $M$ پیش از انجام
$n$ گام متوقف نشده باشد، توصیف~$!C(M,w,n)$ از پیکربندی~$M$ در
زمان~$n$ از~$!T(M,w)$ نتیجه می‌شود. فرض کنید $M$
پس از~$k$ گام متوقف شود. در آن نقطه، برای بعضی~$m\in\Nat$، خانهٔ~$m$
را پویش می‌کند. آنگاه $!C(M,w,k)$ پیکربندی توقف~$M$ را توصیف می‌کند؛
یعنی هم $\Obj Q_q(\num{m},\num{k})$ و هم
$\Obj S_\sigma(\num{m},\num{k})$ را با $\delta(q,\sigma)$ تعریف‌نشده
به‌عنوان عطف در بر دارد. بنابراین بنا بر
\olref{lem:halt-config-implies-halt}،
$!C(M,w,k)\Entails !E(M,w)$. اما چون
$!T(M,w)\Entails !C(M,w,k)$، داریم
$!T(M,w)\Entails !E(M,w)$ و در نتیجه
$!T(M,w)\lif !E(M,w)$ معتبر است.
\end{proof}


\begin{explain}
برای تکمیل بررسی ادعایمان باید جهت عکس را نیز ثابت کنیم: اگر
$!T(M,w)\lif !E(M,w)$ معتبر باشد، $M$ در واقع پس از آغاز با ورودی~$w$
متوقف می‌شود.
\end{explain}

\begin{lem}
\ollabel{lem:halt-if-valid}
اگر $\Entails !T(M,w)\lif !E(M,w)$، آنگاه $M$ روی ورودی~$w$ متوقف
می‌شود.
\end{lem}

\begin{proof}
!!{structure}~$\Struct M$ برای زبان~$\Lang L_M$ را در نظر بگیرید که دامنه‌اش
$\Nat$ است، $\Obj 0$ را به~$0$، $\prime$ را به تابع جانشین و~$<$ را
به رابطهٔ کوچک‌تر تعبیر می‌کند و محمول‌های~$\Obj Q_q$ و
$\Obj S_\sigma$ را چنین تعبیر می‌کند:
\begin{align*}
  \Assign{\Obj Q_q}{M} & =
\Setabs{\tuple{m, n}}{\begin{array}{ll}\text{$M$ پس از آغاز روی~$w$ و اجرای~$n$ گام}\ \text{در حالت $q$ است
  و خانهٔ~$m$ را پویش می‌کند}
\end{array}} \\
\Assign{\Obj S_\sigma}{M} & = \Setabs{\tuple{m, n}}{\begin{array}{ll}
\text{$M$ پس از آغاز روی~$w$ و اجرای~$n$ گام،}\ \text{در خانهٔ~$m$
  نماد~$\sigma$ را دارد}
\end{array}}
\end{align*}
به بیان دیگر، !!{structure}~$\Struct M$ را چنان می‌سازیم که آنچه را $M$
پس از آغاز روی ورودی~$w$ واقعاً گام‌به‌گام انجام می‌دهد توصیف کند.
روشن است که $\Sat{M}{!T(M,w)}$. اگر
$\Entails !T(M,w)\lif !E(M,w)$، آنگاه $\Sat{M}{!E(M,w)}$ نیز؛ یعنی
\[
\Sat{M}{\lexists[x][\lexists[y][(\bigvee_{\tuple{q, \sigma} \in
      X}(\Obj Q_q(x, y) \land \Obj S_\sigma(x, y)))]]}.
\]
چون $\Domain{M}=\Nat$، باید $m,n\in\Nat$ای وجود داشته باشند که برای
بعضی~$q$ و~$\sigma$ با $\delta(q,\sigma)$ تعریف‌نشده،
$\Sat{M}{\Obj Q_q(\num{m},\num{n})\land
\Obj S_\sigma(\num{m},\num{n})}$. بنا بر تعریف~$\Struct M$، یعنی
$M$ پس از آغاز روی ورودی~$w$ و اجرای~$n$ گام در حالت~$q$ است و
نماد~$\sigma$ را می‌خواند، در حالی که تابع گذار تعریف‌نشده است؛ یعنی
$M$ متوقف شده است.
\end{proof}

\end{document}
